\chapter{Boolean Analysis — Circuit Reindex}
%%% generated by the blueprint pipeline from TCSlib.BooleanAnalysis.LMN.CircuitReindex
\section{Overview}

This module provides infrastructure for re-indexing the gate inputs of a Boolean
circuit: it defines an operation that maps the variable indices of a circuit through
a function $f : \mathrm{Fin}\,m \to \mathrm{Fin}\,m'$, and proves that this operation
preserves circuit depth and commutes with evaluation.

%%%%%%%%%%%%%%%%%%%%%%%%%%%%%%%%%%%%%%%%%%%%%%%%%%%%%%%%%%%%%%%%%%%%%%%%%%%%%%%
\section{Declarations}

\begin{definition}[Re-indexing a circuit]
\lean{BoolCircuit.Circuit.reidx}
\label{BoolCircuit.Circuit.reidx}
\leanok
\uses{BoolCircuit.Circuit}
Given a circuit $c$ on $m$ variables and a function $f : \mathrm{Fin}\,m \to \mathrm{Fin}\,m'$,
the re-indexed circuit $\mathrm{reidx}(c, f)$ on $m'$ variables is obtained by replacing
each literal index $i$ by $f(i)$ (keeping its sign) and recursing structurally through the
$\mathrm{And}/\mathrm{Or}$ nodes.
\end{definition}

\begin{theorem}[Re-indexing preserves circuit depth]
\lean{BoolCircuit.Circuit.reidx_depth}
\label{BoolCircuit.Circuit.reidx_depth}
\leanok
\uses{BoolCircuit.Circuit, BoolCircuit.Circuit.depth, BoolCircuit.Circuit.ind, BoolCircuit.Circuit.reidx}
\difficulty{5}
Let $c$ be a Boolean circuit on $m$ variables and let $f : \mathrm{Fin}\,m \to
\mathrm{Fin}\,m'$ be any map of variable indices. Then re-indexing $c$ along $f$ leaves
its depth unchanged: $\mathrm{depth}(\mathrm{reidx}(c, f)) = \mathrm{depth}(c)$.
\end{theorem}

\begin{theorem}[Re-indexing commutes with evaluation]
\lean{BoolCircuit.Circuit.reidx_eval}
\label{BoolCircuit.Circuit.reidx_eval}
\leanok
\uses{BoolCircuit.Circuit, BoolCircuit.Circuit.eval, BoolCircuit.Circuit.ind, BoolCircuit.Circuit.reidx, BoolCircuit.Lit.eval}
\difficulty{5}
Let $c$ be a Boolean circuit on $m$ variables, let $f : \mathrm{Fin}\,m \to
\mathrm{Fin}\,m'$ be a re-indexing map, and let $g : \mathrm{Fin}\,m' \to \bbf_2$ be an
assignment to the $m'$ variables. Then evaluating the re-indexed circuit
$\mathrm{reidx}(c, f)$ under $g$ gives the same value as evaluating the original circuit
$c$ under the pulled-back assignment $g \circ f$:
\[
\mathrm{eval}\big(\mathrm{reidx}(c, f),\, g\big) = \mathrm{eval}\big(c,\, g \circ
f\big).
\]
\end{theorem}
